\pg{430} a) As substancias creadas estão sujeitas a continuos movimentos ou mudanças. Ora os movimentos, sendo passagens da \emph{potencia} para o \emph{acto}, devem ser determinados --- não pelas proprias substancias que se movem, pois ninguem dá o que não tem, --- mas por um motor distincto, pois o ente, dotado apenas de \emph{potencia}, não adquire o \emph{acto}, se lh'o não communica um outro ente, que já o possue. Este motor --- ou é \emph{movido}, --- ou é \emph{immovel}. Se é \emph{immovel}, é \textsc{Deus}. Se é \emph{movido}, subiremos de motor em motor, até encontrarmos necessariamente um motor \emph{immovel}; porque, se não existisse este motor \emph{immovel}, não poderiam existir os motores \emph{moveis}, que não têm em si a razão do seu movimento, e por isso não haveria movimento nas substancias creadas. Logo existe \textsc{Deus} \footnote{Aqui a palavra \emph{movimento} toma-se em sentido mais lato, emquanto exprime não só o \emph{movimento local}, mas toda e qualquer \emph{mudança} ou \emph{passagem da potencia para o acto}. D'este modo o argumento exposto equivale a este: O mundo é \emph{mutavel}. Ora todo o ente \emph{mutavel} depende do ente \emph{immutavel}. Logo existe o ente \emph{immutavel}. O ente \emph{immutavel} é \textsc{Deus}. Logo existe \textsc{Deus}. --- Ora, encontrando-se a \emph{mudança} não só nos corpos mas tambem nos espiritos, o argumento póde desdobrar-se; pois, baseando-nos tanto no movimento dos \emph{corpos} como no dos \emph{espiritos}, podemos elevar-nos ao descobrimento de uma causa, que move tudo, sendo ella \emph{immovel} ou \emph{immutavel}.\par Eis os dois argumentos:\par a) A experiencia attesta que as substancias materiaes, por exemplo, os planetas, estão sujeitas a continuos movimentos. Ora todo o corpo que se move, é movido por um agente; pois o corpo, sendo por si indifferente tanto para a quietação como para o movimento, não póde determinar-se por si ao movimento. Logo ou devemos admittir uma serie infinita de motores \emph{moveis}, ou deve existir um motor \emph{immovel}, que seja a causa de todo o movimento dos corpos. Ora uma serie infinita de motores \emph{moveis} repugna; porque numa serie \emph{infinita} de motores \emph{moveis} não haveria um \emph{primeiro} motor, e assim não existiria uma causa proporcionada ou adequada de todos os movimentos\pgn{431}nos corpos. Logo deve existir um motor \emph{immovel}. Este motor \emph{immovel} deve ser infinitamente perfeito; pois, se o não fosse, se lhe faltasse alguma perfeição, poderia tender para ella, e assim passaria da potencia ao acto, e não seria \emph{immovel}. O ser infinitamente perfeito é \textsc{Deus}. Logo existe \textsc{Deus}.\par Este argumento, muito popular, pode desenvolver-se \emph{scientificamente} da maneira seguinte. Existe no mundo o movimento. Para a explicação da existencia do movimento adduzem-se duas hypotheses: a \emph{mechanica} e a \emph{dynamica}. Ora nenhuma d'estas hypotheses pode, na explicação do movimento, prescindir da intervenção de um primeiro Motor. --- Não pode prescindir a hypothese \emph{mechanica}. Na verdade, para os sequazes d'esta hypothese, a materia é \emph{meramente passiva}, de modo que os corpos são movidos por um impulso extrinseco, mas elles não movem, e só communicam aos outros corpos o movimento recebido. Se todos os corpos, de que se compõe o universo, são meramente passivos, é necessario que a força impulsiva do movimento venha de fóra, e proceda de um primeiro Motor, que não seja movido. Nem se diga que o movimento da materia teve um começo \emph{necessario}, ou que é \emph{eterno}; porque um \emph{começo necessario} é uma contradicção (o que começa é necessariamente contingente), e um movimento \emph{eterno} da materia, embora não envolvesse repugnancia, não deixaria de exigir um primeiro motor, desde que se supponha a materia meramente \emph{passiva}. --- Não pode prescindir a hypothese \emph{dynamica}. Com effeito, para os sequazes d'esta hypothese, a materia não é meramente passiva, mas é essencialmente dotada de actividade, em virtude da qual um corpo pode, em certas circumstancias, mover outro corpo, \emph{attrahendo-o}. E' certo que o mundo inteiro se move pela \emph{lei da attracção}; mas a \emph{attracção} não é o \emph{primeiro} motor, nem é o motor \emph{sufficiente} do nosso systema solar. Não é o \emph{primeiro} motor; porquanto, a força de \emph{attracção}, de que é dotada a materia, não está sempre e necessariamente \emph{em acto}, mas, pelo contrario, tende para o repouso das partes numa unica massa, e por isso é indifferente para o acto e para a potencia, e, não conténdo em si a razão da passagem da potencia para o acto, exige e suppõe uma causa que a determine e mova. Não é o motor \emph{sufficiente} do nosso systema solar; porquanto este systema deve ser mantido em equilibrio e regulado nos seus movimentos por duas forças --- uma \emph{centripeta} (sem a qual não haveria unidade no systema) e outra \emph{centrifuga} (sem a qual todos os corpos confundir-se-hiam numa unica massa), e se a força \emph{centripeta}, da qual deriva a attracção, pode dizer-se essencial á materia, não se pode dizer essencial a força \emph{centrifuga}, a qual por isso deve ser accrescentada e unida por um motor extranho e superior á propria materia.\par Dizem que este argumento, se prova a existencia de um Ente superior ao mundo e causa de todo o movimento, não prova a existencia de um Deus tão perfeito, como o concebem os espiritualistas. --- Respondemos que o nosso argumento, embora não provasse senão a \emph{existencia} de um Ser supremo, já teria alcançado o seu fim, que é o de demonstrar a mesma \emph{existencia}, e não a \emph{natureza} do Ente supremo. São duas questões distinctas. --- Todavia devemos admittir que este argumento, que prova a necessidade de um \emph{primeiro Motor}, contém, \emph{em germe}, a razão de todos os attributos da Divindade.\pgn{432}Na verdade, se a materia, por ser indifferente ao repouso e ao movimento, tem a razão do seu movimento não em si mas numa causa exterior, segue-se que tal causa existe fóra da materia, e por isso é immaterial, intelligente, eterna, necessaria, e tambem infinita e unica. Porquanto, o primeiro Motor, não podendo ser movido, é \emph{todo acto}, sem mistura de potencialidade, é acto \emph{purissimo}, isto é, infinitamente perfeito e é unico.\par Os melhores naturalistas e philosophos modernos reconhecem esta verdade. Citaremos apenas \emph{Couchy}, \emph{Clarke}, \emph{B. S. Hilaire}, \emph{P. Janet}. No fim, ouviremos um insuspeito, \emph{Rousseau}.\par \emph{Couchy} diz: « O unico ser, que póde produzir as forças physicas, é o Ente necessario. A força é a expressão da sua vontade. As diversas forças do equilibrio e movimento são causas secundarias. A gravitação universal, a attracção dos corpos na razão inversa do quadrato das distancias, que os separam, é uma lei que exprime a vontade divina » (\emph{Comptes rendus hebd. de l'Acad. des sciences}, t. XXI, Paris, 1845).\par \emph{Clarke}, discipulo de \emph{Newton}, fallando da \emph{attracção}, de que depende todo o movimento dos corpos celestes, escreve: « Ella póde ser o effeito de um impulso, mas não póde ser material, é antes o effeito de uma causa immaterial » (\emph{Physica de Rohault}, trad. por \emph{Clarke}, II, c. XI).\par \emph{Barthélemy Saint-Hilaire} diz: « Concebe-se que a materia siga hoje as leis infalliveis, que a regem, e conserve uma regularidade eterna; mas foi necessario um impulso primitivo, que ordenou tudo para a inexhaurivel serie dos tempos. Até as anomalias, que offerece o admiravel systema dos ceus, attestam a presença immanente e indefectivel d'Aquelle que os fez. O sobrenatural está em tudo... Só é necessario que a sciencia se resigne a resolver certos problemas por outros meios, que não sejam os de uma experiencia impossivel; e o da origem de todas as coisas é um d'esses problemas, a que só se renuncia por timidez, pensando que se pratica uma reserva prudente. \emph{A questão da origem é inevitavel}, e de nada valeria querer illudil-a ». (\emph{Journal des Savants}, 1862, pag. 607).\par \emph{P. Janet} diz: « A experiencia não mostra senão que um corpo é movido por outro, e este por outro, e isto ao infinito. Ora, se fóra d'esta cadêa não existe uma causa motora, teremos uma serie de movimentos sem causa e sem razão » (\emph{Le matérialisme contemporain}, 1875, pag. 58).\par \emph{Rousseau} escreve: « Vejo a materia ora em movimento e ora em repouso; d'ahi concluo que nem o repouso nem o movimento lhe são essenciaes. Mas o movimento, sendo uma acção, mostra ser effeito de uma causa... Conceber a materia productora do movimento é conceber um effeito destituido de causa... Não é evidente que, se o movimento fosse essencial á materia, um não poderia separar-se da outra, mas estaria sempre unido com ella, conservando-se sempre no mesmo grau, sempre o mesmo em cada uma das partes, não poderia augmentar nem diminuir, e não se poderia conceber a materia em repouso?!... Quando me dizem que o movimento não é \emph{essencial} á materia, mas \emph{necessario}, pretendem embaraçar-me com palavras, que poderiam facilmente confutar-se, se tivessem algum sentido. Porque o movimento da materia --- ou deriva da propria materia, e então lhe é essencial\pgn{433} --- ou deriva de uma causa extrinseca, e já não é necessario á propria materia, senão emquanto a causa motora actua nella » (\emph{Emilio}, t. III).\par Na \emph{Cosmologia geral} tambem desenvolvemos este ponto.\par b) No mundo não só existe o movimento \emph{local}, mas tambem o movimento \emph{espiritual}; e este, como aquelle, suppõe e exige necessariamente a existencia de um primeiro motor immovel. Na verdade, se reflectirmos em nós mesmos, percebemos que a nossa alma tambem se move, emquanto passa continuamente da potencia para o acto. O acto é um complemento, uma addição que aperfeiçõa a potencia. Ora este acto não pode ser e não é extrahido do fundo da potencia, porque ninguem dá o que não tem; mas deve provir de um agente extrinseco, de um motor, que possua \emph{actualmente} o que falta á potencia e que seja capaz de lh'o communicar. Este motor é a propria alma, que se move a si mesma emquanto uma faculdade, que está \emph{em acto}, move uma outra, que está \emph{em potencia}. Todavia a alma não pode ser o \emph{primeiro} motor. O primeiro acto da faculdade, que move todas as outras, não pode ser o effeito do impulso das outras faculdades, que estão em potencia, mas deve ser o effeito de um motor distincto e superior á propria faculdade. Este motor --- ou é o primeiro Motor, --- ou é um motor intermedio, movido, finalmente, por esse primeiro Motor immovel.\par Este argumento, que se refere ao movimento da alma humana em geral, conserva toda a sua força, quando se refere ao movimento da \emph{intelligencia} e da \emph{vontade}.\par \emph{a}) Dá-se o movimento na \emph{intelligencia}. Ainda que a intelligencia mova a vontade, e esta mova aquella, comtudo, por ser absurdo o processo ao infinito, devemos admittir que o primeiro acto da intelligencia precedeu o primeiro acto da vontade. Diz S. Thomaz: « Non oportet procedere in infinitum, sed statur in intellectu, sicut in primo. Omnem enim voluntatis motum necesse est quod praecedat apprehensio; sed non omnem apprehensionem praecedit motus voluntatis » (\emph{Sum. Th.}, p. I, q. 82, a. 4, ad 3). Mas de que modo passa a intelligencia da potencia para o acto? Espontaneamente? Não; porque faltaria a razão sufficiente da passagem. Pela acção do objecto intelligivel? Tambem não; porque, sendo este objecto formado pela intelligencia agente, a acção do objecto é \emph{posterior} á da intelligencia. A causa motora deve ser superior ao objecto e a sujeito, ao intelligivel e ao intelligente; e esta causa é Deus. Aristoteles, ainda que pagão, exprime esta verdade nas seguintes magnificas palavras: « Quod est principium motus in anima? Patet vero, quemadmodum in universo est Deus, sic etiam in anima; moventur enim omnia ab inexistente in nobis Numine. Intellectus autem principium non intellectus sed praestantius quidpiam est. Quid autem et scientia praestantius est et ratione, nisi Deus? (\emph{Moral. Eudem.}, l. 7, c. 14). D'onde se vê a verdade do aphorismo: \emph{Eu penso; logo Deus existe.}\par \emph{b}) Não só o acto da intelligencia, mas tambem o da \emph{vontade} exige e suppõe a existencia de um primeiro Motor. A vontade é movida necessariamente pelo bem absoluto. Esse movimento provoca uma tendencia \emph{universal}, \emph{necessaria}, \emph{indestructivel}. Ora um effeito, dotado de semelhantes qualidades, como é essa tendencia, não pode derivar de uma causa imaginaria ou phantastica, mas deve derivar de uma causa real e objectiva, que deve\pgn{434}identificar-se com o proprio bem absoluto; porque todo o effeito suppõe uma causa proporcionada. Ora o bem absoluto é Deus. Logo Deus existe.\par Como se vê, o principio --- \emph{tudo o que se move é movido por um motor distincto e, em ultima analyse, immovel} --- é muito fecundo e serve para a intelligencia e resolução de innumeraveis problemas.}.

b) No mundo existem muitas e diversas \emph{causas efficientes}. Estas causas estão por tal modo subordinadas entre si, que formam uma serie, na qual uma depende de outra, a inferior \pg{431} \pg{432} é \emph{effeito} da superior. Ora, nesta serie de causas \emph{efficientes}, deve existir uma causa \emph{primeira}, que seja só \emph{causa} e não \emph{effeito}, que seja causa de todas as causas; visto que, se não existisse uma \pg{433} causa \emph{primeira}, não poderiam existir, como existem, as causas \emph{segundas}, pois as causas \emph{segundas} são \emph{effeitos}, e todo o effeito